%%%%%%%%%%%%%%%%%%%%%%%%%%%%%%%%%%%%%%%%%%%%%%%%%%%%%%%%%%%%%%%%%%%%%%%%%%%%%%%%%%%%%%%%%%%%%%%%%%%%%%%%
%                                                                                                      %
%% Modelo desenvolvido pelos Alunos e Professores do IFSP - PEP, para                                  %
% MOSTRA CIENTÍFICA, CULTURAL E TECNOLÓGICA 2021.                                                      % 
%                                                                                                      % 
% Alunos: Christopher Alves de Oliveira, Inácio Ribeiro, Mateus Antônio Guelfi.                        %    
% Professores: Andryos da Silva Lemes, Leonardo Ataide Carniato.                                       % 
% Versão: 1.0                                                                                          % 
% Data: 27/08/2020.                                                                                    %
%                                                                                                      %
%%%%%%%%%%%%%%%%%%%%%%%%%%%%%%%%%%%%%%%%%%%%%%%%%%%%%%%%%%%%%%%%%%%%%%%%%%%%%%%%%%%%%%%%%%%%%%%%%%%%%%%%

% Selecionando a classe que gera artigos, tamanho da folha, tamanho padrão letras, coluna dupla e etc.
\documentclass[a4paper,10pt,twocolumn,pdftex]{article} 

% Pacote utilizado para definir tamanho das seções, subseções, adicionar linhas e etc
\usepackage{titlesec} 
\titleformat{\section}[block]{\centering\normalsize\bfseries}{\thesection}{1pt}{}[\hrule]
\titleformat{\subsection}{\centering\normalsize\bfseries}{\thesubsection}{1pt}{}
\titleformat{\subsubsection}{\centering\normalsize\bfseries}{\thesubsubsection}{1pt}{}

\usepackage{array}
\usepackage{booktabs}

% Inclui as bibliotecas principais
\input{Layout/preambulo} 

% Comando utilizado para facilitar a inserção do texto "Fonte"
\newcommand{\fonte}{\footnotesize\textit{Fonte:}} 

% Comando utilizado para facilitar a inserção do texto "Elaboração do autor" 
\newcommand{\elaborado}{\footnotesize Elaborada pelo autor.} 

 % Inicia o documento
\begin{document}
\pagestyle{empty} %Iniciar o documento vazio.

%----------------------------------------TEXTOS------------------------------------------------------------

% Comando que gera colunas falsas no cabeçalho para não dividir e sobrepor os textos
\twocolumn[\begin{@twocolumnfalse} 

\input{Layout/cabeçalho} % Alterar a área de interesse, inserir/retirar logo da instituição e alterar Título. 

\input{Layout/dados} % Alterar nomes e e-mails dos autores. OBS: Comentar essa linha antes de enviar o artigo para avaliação.

% Comando que gera colunas falsas no cabeçalho para não dividir e sobrepor os textos
\end{@twocolumnfalse}] 

\input{Textos/textos} % Alterar, resumo, metodologia, resultados e agradecimentos.

%-----------------------------------------------------------------------------------------------------------


%-----------------------------------REFERÊNCIAS-------------------------------------------------------------

{\bibliography{bibliografia}} % Biblioteca responsável por gerir as referências, os autores recomendam utilizar em conjunto com o software Mendeley. 

%----------------------------------------------------------------------------------------------------------

 % Finaliza o documento
\end{document}
